% =====================================================================
%  CROSS-DOMAIN WI-FI CSI PERCEPTION VIA SELF-SUPERVISED DISENTANGLEMENT
%  AND SPATIAL RETRIEVAL-AUGMENTED GENERATION
%
%  Target: IEEE conference paper (IEEEtran, two-column, 6-8 pages incl. refs)
%
%  ---------------------------------------------------------------------
%  !!  READ BEFORE SUBMITTING  !!
%  This draft contains NO measured numbers. Every place where a result
%  belongs is marked with a red [RESULT NEEDED: ...] placeholder, and
%  every experimental design choice is marked [DECISION]. Nothing in this
%  file was invented: the numbers must be filled from real runs, and the
%  citations were fetched programmatically into references_verified.bib.
%  Set \draftnoticefalse at the end to remove the draft banner.
%  ---------------------------------------------------------------------
% =====================================================================

\documentclass[conference]{IEEEtran}

% ---- packages (kept minimal to avoid IEEEtran conflicts) -------------
\usepackage{cite}
\usepackage{amsmath,amssymb}
\usepackage{graphicx}
\usepackage{tikz}
\usepackage{booktabs}
\usepackage{xcolor}
\usepackage{url}
\usepackage{algorithm}
\usepackage{algpseudocode}
% \usepackage{subcaption}  % removed: unused, and not available in the CTAN tlnet archive
% \usepackage{balance}  % removed: not available in the CTAN tlnet archive

% ---- draft guard -----------------------------------------------------
\newif\ifdraftnotice
\draftnoticetrue

% ---- method name + placeholder macros --------------------------------
\newcommand{\method}{\textsc{Relay}}
\newcommand{\RESULTNEEDED}[1]{\textcolor{red}{\textbf{[RESULT NEEDED: #1]}}}
\newcommand{\DECISION}[1]{\textcolor{blue}{\textbf{[DECISION: #1]}}}
\newcommand{\TODO}[1]{\textcolor{orange}{\textbf{[TODO: #1]}}}

% ---- math shorthands -------------------------------------------------
\newcommand{\zc}{\mathbf{z}^{c}}      % content (gesture/activity) factor
\newcommand{\zd}{\mathbf{z}^{d}}      % domain (environment/layout) factor
\newcommand{\bs}[1]{\boldsymbol{#1}}

\begin{document}

% Cover page, typeset in LaTeX. The manuscript below is unmodified.
\onecolumn
\thispagestyle{empty}
\null
\begin{tikzpicture}[remember picture,overlay]
  \coordinate (covercenter) at (current page.north);
  \node[anchor=north,font=\bfseries\fontsize{10}{12}\selectfont]
    at ([yshift=-0.10in]covercenter) {PROPOSAL SKRIPSI ARTIKEL ILMIAH};

  \node[anchor=north,align=center,text width=6.6in,font=\bfseries\fontsize{9}{13}\selectfont]
    at ([yshift=-0.87in]covercenter)
    {Persepsi CSI Wi-Fi Lintas Domain melalui Pemisahan Representasi Mandiri dan \textit{Spatial}};
  \node[anchor=north,font=\bfseries\itshape\fontsize{9}{12}\selectfont]
    at ([yshift=-1.28in]covercenter) {Retrieval-Augmented Generation};

  \node[anchor=north,align=center,text width=6.4in,font=\bfseries\fontsize{9}{14}\selectfont]
    at ([yshift=-2.06in]covercenter)
    {Cross-Domain Wi-Fi CSI Perception via Self-Supervised Disentanglement and Spatial Retrieval-};
  \node[anchor=north,font=\bfseries\fontsize{9}{12}\selectfont]
    at ([yshift=-2.47in]covercenter) {Augmented Generation};

  \node[anchor=north,font=\fontsize{9}{12}\selectfont]
    at ([yshift=-3.31in]covercenter) {\textbf{Topik:} Intelegensia Semu};

  \node[anchor=north,font=\fontsize{9}{12}\selectfont]
    at ([yshift=-4.10in]covercenter)
    {2702387973 / CHRISTOPHER RICKY SOEHARTONO / Computer Science};
  \node[anchor=north,font=\fontsize{9}{12}\selectfont]
    at ([yshift=-4.52in]covercenter)
    {2702363651 / FRITZ GRADIYOGA / Computer Science};
  \node[anchor=north,font=\fontsize{9}{12}\selectfont]
    at ([yshift=-4.94in]covercenter)
    {2702363765 / JESSICA NATHANIA WENARDI / Computer Science};

  % BINUS emblem and lettering drawn as editable vector elements.
  \begin{scope}[shift={([yshift=-5.72in]covercenter)},x=1cm,y=1.6cm]
    \draw[gray!65,line width=1.1pt] (-3.05,-0.72) -- (3.05,-0.72);
    \draw[red!85!black,line width=.55pt] (0,0.10) -- (0,-2.32);
    \fill[red!85!black] (0,0.10) circle (1.5pt);
    \fill[red!85!black] (0,-2.32) circle (1.5pt);
    \fill[red!85!black] (-3.05,-0.72) circle (1.5pt);
    \fill[red!85!black] (3.05,-0.72) circle (1.5pt);
    \draw[gray!70,line width=3.0pt] (0.15,-.34) arc[start angle=65,end angle=305,x radius=.92,y radius=.70];
    \fill[orange!90!yellow] (0,-.76) circle (.40cm);
    \fill[white] (-.15,-.57) circle (.055cm);
    \foreach \x/\y in {.55/-.20,.72/-.14,.81/-.27,.98/-.13,1.10/-.42,
       .74/-.45,.90/-.55,1.12/-.62,.62/-.72,.81/-.81,1.23/-.30,
       .99/-.77,1.34/-.11,1.52/-.50,1.38/-.76,1.68/-.19}
       \fill[gray!60] (\x,\y) rectangle ++(.055,.055);
    \node[anchor=west,font=\sffamily\bfseries\fontsize{24}{26}\selectfont,text=black!88]
       at (.45,-1.65) {BINUS};
    \node[anchor=west,font=\sffamily\bfseries\fontsize{19}{21}\selectfont,text=cyan!75!blue]
       at (.47,-2.68) {UNIVERSITY};
  \end{scope}

  \node[anchor=north,font=\bfseries\fontsize{9}{12}\selectfont]
    at ([yshift=-8.44in]covercenter) {BINUS University};
  \node[anchor=north,font=\bfseries\fontsize{9}{12}\selectfont]
    at ([yshift=-8.87in]covercenter) {2026};
  \node[anchor=north west,font=\fontsize{9}{12}\selectfont]
    at ([xshift=.98in,yshift=-9.51in]current page.north west) {DIPERIKSA OLEH};
  \node[anchor=north west,align=left,font=\fontsize{8}{9}\selectfont]
    at ([xshift=.73in,yshift=-10.61in]current page.north west)
    {D3557 - Ir. Andry Chowanda, S.Kom, MM,\\
     Ph.D., MBCS, CCP, CME, IPM, SMIEEE};
\end{tikzpicture}
\clearpage
\twocolumn


\title{Cross-Domain Wi-Fi CSI Perception via Self-Supervised\\
Disentanglement and Spatial Retrieval-Augmented Generation}

\author{
\begin{tabular}[t]{ccc}

\begin{tabular}[t]{c}
\IEEEauthorblockN{Christopher Ricky Soehartono} \\[-0.8cm]
\IEEEauthorblockA{\textit{Computer Science Department} \\
\textit{School of Computer Science} \\
\textit{Bina Nusantara University} \\
Jakarta, Indonesia 11480 \\
christopher.soehartono@binus.ac.id}
\end{tabular}

&

\begin{tabular}[t]{c}
\IEEEauthorblockN{Fritz Gradiyoga} \\[-0.8cm]
\IEEEauthorblockA{\textit{Computer Science Department} \\
\textit{School of Computer Science} \\
\textit{Bina Nusantara University} \\
Jakarta, Indonesia 11480 \\
fritz.gradiyoga@binus.ac.id}
\end{tabular}

&

\begin{tabular}[t]{c}
\IEEEauthorblockN{Jessica Nathania Wenardi} \\[-0.8cm]
\IEEEauthorblockA{\textit{Computer Science Department} \\
\textit{School of Computer Science} \\
\textit{Bina Nusantara University} \\
Jakarta, Indonesia 11480 \\
jessica.wenardi@binus.ac.id}
\end{tabular}

\end{tabular}

\\[1cm]

\begin{tabular}[t]{c}
\IEEEauthorblockN{Ir. Andry Chowanda, S.Kom, MM, Ph.D.,} \\[-0.8cm]
\IEEEauthorblockN{MBCS, CCP, CME, IPM, SMIEEE} \\[-0.8cm]
\IEEEauthorblockA{\textit{Computer Science Department} \\
\textit{School of Computer Science} \\
\textit{Bina Nusantara University} \\
Jakarta, Indonesia 11480 \\
achowanda@binus.edu}
\end{tabular}
}

\maketitle

\ifdraftnotice
\begin{center}
\fbox{\parbox{0.94\columnwidth}{\small\textbf{\textcolor{red}{DRAFT --- RESULTS NOT YET MEASURED.}}
This manuscript states the method, the evaluation protocol, and the positioning against
prior work. It reports \emph{no} empirical numbers, because no experiments have been run
yet. All result slots are marked \RESULTNEEDED{...}. See \texttt{HANDOFF.md} for the
runnable experiment plan.}}
\end{center}
\vspace{-0.4em}
\fi

% =====================================================================
\begin{abstract}
% 5-sentence formula (Farquhar): achieved / why hard / how / evidence / headline number.
We introduce \method{}, a cross-domain Wi-Fi channel state information (CSI) sensing
framework that couples \emph{self-supervised disentanglement} of gesture content from
environmental conditions with a \emph{spatial retrieval-augmented} memory that supplies
the model with an explicit, updatable description of the space it currently occupies.
Existing CSI sensing systems are trained and deployed in the same room: because the
received signal is a function of both the human-induced dynamics and the static
propagation environment, accuracy collapses when the layout, transmitter placement, or
furniture configuration changes, and current remedies either require labelled target data,
a retraining round, or multiple transceivers.
Our key idea is to treat the environment not as nuisance to be erased but as
\emph{retrievable context}: a self-supervised objective separates a content code from a
domain code, and at inference a geometry-tagged memory indexed by coarse spatial
descriptors conditions the classifier on the active layout, so a new room becomes a
retrieval problem rather than a retraining problem.
We evaluate on \DECISION{dataset(s) to be finalised --- candidates: Widar3.0, SignFi, Wiar,
CSIDA, XRF55} under cross-environment, cross-location, and cross-orientation protocols,
against domain-adaptation and domain-generalisation baselines.
\RESULTNEEDED{abstract headline: best cross-domain accuracy/F1 with mean\,$\pm$\,std over
$N$ seeds, and the absolute gain over the strongest baseline}.
\end{abstract}

\begin{IEEEkeywords}
Wi-Fi sensing, channel state information, cross-domain generalisation, self-supervised
learning, disentangled representation, retrieval-augmented generation.
\end{IEEEkeywords}

% =====================================================================
\section{Introduction}
\label{sec:intro}

Wi-Fi sensing infers human activity from the channel state information (CSI) that
commodity network interface cards already expose, and it is attractive precisely because
it needs no new hardware, works through walls, and observes no images
\cite{Chen2022Survey,Radwan2025Survey}. Its central weakness is equally well documented:
a CSI sample mixes what the person is doing with where the measurement is taken. The same
gesture performed in a different room, at a different distance from the transmitter, or
after a door is opened produces a CSI trajectory that a classifier trained in the original
setting misreads. Reported cross-domain drops are severe --- the survey of Chen
et al.\ \cite{Chen2022Survey} and the more recent taxonomy of Wang et al.
\cite{Wang2025GeneralizabilitySurvey} both organise an entire literature around the
problem of maintaining accuracy when "the deployment domain is not the training domain",
and Zhang et al.\ \cite{Zhang2026BeyondLabels} argue the very definition of a domain
under-specifies the factors that matter.

Prior work attacks this along a well-charted axis. \emph{Domain adaptation} aligns source
and target feature distributions from unlabelled target data \cite{Yang2021BookChapter,
Chen2022Fidora,Mao2024WiCro,Prasad2024DANN}, \emph{domain generalisation} instead seeks
invariance from source data alone \cite{Liu2024WiSR,Liu2024WiSGP,Wang2022AirFi,
Xu2023ARCFi}, \emph{few-shot} methods match a handful of target samples against learned
prototypes \cite{Yin2022FewSense,Zhao2024CrossFi,Yang2023FewCS}, and \emph{self-supervised}
pretraining turns cheap unlabelled CSI into a transferable encoder
\cite{Yang2022AutoFi,Barahimi2024CPC,Lyons2024WiFiAct,Radwan2025Survey}. Two observations
motivate this paper.

First, invariance is a lossy target. When a model is trained to be blind to the
environment, it discards information that is genuinely useful for deciding whether a
reading is plausible --- a fall signature in a narrow corridor and in an open hall are not
the same measurement, and a system that cannot tell the two apart has no way to notice
that it is being asked about a space it has never modelled. Disentanglement decomposes
this problem rather than deleting it \cite{Kang2021ContextAware}, but the domain factor is
then usually thrown away, used only adversarially.

Second, the field already names the missing ingredient. The generalisability survey of
Wang et al.\ \cite{Wang2025GeneralizabilitySurvey} explicitly asks whether
retrieval-augmented generation (RAG) can act as "an external database" that "allows the
model to access Wi-Fi knowledge on demand". In the language-model literature that idea is
mature \cite{lewis2020rag}, and it has begun to reach wireless systems
\cite{Nazar2024Enwar,Nazar2026Enwar2,wi_chat} and spatially-constrained reasoning
\cite{spatial_rag}. To our knowledge no work uses retrieval over \emph{spatial/layout}
descriptions as the adaptation mechanism for cross-domain CSI \emph{sensing}: retrieval has
been applied to network management and to interpreting CSI with language models, not to
supplying the sensing head with the layout of the room it is currently in.

We therefore propose \method{}, which (i) learns a self-supervised factorisation of CSI
into a \emph{content} code carrying gesture/activity information and a \emph{domain} code
carrying the layout and hardware conditions, and (ii) at inference retrieves, from a
geometry-tagged memory, the layout entries most consistent with the current signal
statistics and conditions the classifier on them. Adaptation to a new space then reduces
to inserting that space into the memory once --- no labelled target data and no gradient
retraining --- which is the regime that deployment actually imposes.

\subsection*{Contributions}
\begin{itemize}
\item We formulate cross-domain CSI sensing as \emph{retrieval-conditioned
perception}: the domain code is preserved as a retrieval key instead of being discarded,
and a new environment is handled by memory insertion rather than retraining
(Section~\ref{sec:method}).
\item We give a self-supervised disentanglement objective for CSI with an explicit
independence penalty and physics-informed augmentations, and we state precisely which
invariance each term buys (Section~\ref{sec:method}).
\item We specify a falsifiable evaluation protocol with claim-to-experiment mapping,
strong DA/DG/few-shot baselines, and a memory-free ablation that isolates the retrieval
contribution (Section~\ref{sec:protocol}), and we report \DECISION{statistical tests to be
fixed: paired bootstrap / McNemar} rather than single-run accuracies.
\end{itemize}

\TODO{contribution 3 currently promises an evaluation that has not been executed; either
run it or soften the claim before submission.}

% =====================================================================
\section{Related Work}
\label{sec:related}

\subsection{Cross-domain Wi-Fi sensing: the problem and its taxonomy}
CSI is a sampled complex channel response; its magnitude and phase respond both to human
motion and to the static geometry of the deployment. Chen et al.\ \cite{Chen2022Survey}
formalise the domain shift this creates and group remedies into five families ---
domain-invariant feature extraction, virtual sample generation, transfer learning,
few-shot learning, and big-data solutions --- and find each has distinct limitations.
Wang et al.\ \cite{Wang2025GeneralizabilitySurvey} review over 200 papers by pipeline
stage, covering domain adaptation, meta-learning, metric learning, augmentation,
cross-modal alignment, federated and continual learning, and release a dataset platform;
they list pretraining at scale, multimodal foundation models, and continual deployment as
the open directions. Radwan et al.\ \cite{Radwan2025Survey} survey specifically the
self-supervised variants, contrasting contrastive and non-contrastive objectives, and
publish a comparative accuracy analysis; Sai et al.\ \cite{Sai2026Review} give a broader
review of machine-learning methods for CSI recognition. Zhang et al.\
\cite{Zhang2026BeyondLabels} argue that the conventional factorisation of domains
(room / user / orientation) is itself the limitation, and redefine domains to make
generalisation claims measurable. We adopt that critique: our formulation keeps the
domain as an explicit, retrievable variable rather than treating all shift as one
undifferentiated nuisance.

\subsection{Domain adaptation and generalisation}
Adaptation-based methods assume some target access. Fidora \cite{Chen2022Fidora} combines a
variational-autoencoder augmenter with a joint classification--reconstruction classifier
and reports gains for new users (+17.8\% average F1) and for a changed environment
(beating the prior state of the art by 23.1\%). Yang et al.\ \cite{Yang2021IGLDA} raise
localisation accuracy from 26.3\% to 82.2\% by recollecting 37.5\% of the data under an
improved global/local metric method. CrossSense \cite{Zhang2018CrossSense} generates
synthetic per-environment training samples with a roaming model plus a mixture of experts,
lifting gait and gesture accuracy from 20\% to over 80--90\%. AdaWiFi
\cite{Zheng2025AdaWiFi} lets multiple IoT devices sense collaboratively and reports
roughly 30\% higher accuracy under one-shot adaptation. Prasad et al.\
\cite{Prasad2024DANN} adversarially learn domain-invariant compressed RSS features with a
gradient reversal layer, and Zhang et al.\ \cite{Zhang2025OnlineDFL} maintain an updatable
radio map with a forgetting mechanism so that sequential data from other times and
layouts can be absorbed without retraining. Wi-Cro \cite{Mao2024WiCro} translates source
data into the target style with a GAN and applies metric learning; Avola et al.\
\cite{Avola2025RaGAN} suppress domain-specific features with a relativistic-average GAN
validated inside a constructed Faraday cage. Source-free adaptation removes the need for
source data: Wi-SFDAGR \cite{Yan2025WiSFDAGR} treats adaptation as unsupervised clustering
with an attraction--dispersion network and uncertainty-weighted neighbours.
Li et al.\ \cite{Li2026Conformal} sidestep retraining with conformal prediction over
multivariate kernel density estimates, reporting accuracy improvements from 30\% to 74\%
across six kinds of domain variation, and He et al.\ \cite{He2025CrossSubject} transfer
across subjects, while Mehryar \cite{Mehryar2023PassiveWiFi} places a few-shot adaptation
stage above a cross-modal recognition layer for deployment-time robustness.
\method{} differs in that it neither requires unlabelled target data at
adaptation time nor retrains: it conditions on a retrieved description of the space.

Generalisation-based methods need only source data. Widar3.0 \cite{Zheng2019Widar3}
estimates body-coordinate velocity profiles as a physically domain-independent
representation and reaches 82.6--92.4\% cross-domain with one-time training. WiSR
\cite{Liu2024WiSR} randomises subcarrier-domain "styles" adversarially, and WiSGP
\cite{Liu2024WiSGP} perturbs inputs along subdomain directions, both targeting
multi-factor generalisation without target data. AirFi \cite{Wang2022AirFi} learns the
critical CSI part for unseen environments; ARC-Fi \cite{Xu2023ARCFi} targets the more
realistic semi-supervised setting where unlabelled CSI is plentiful but labels are scarce,
and blocks "shortcut learning" with an antenna-response-consistency augmentation that
treats co-located antennas as natural views. XFall \cite{Chi2024XFall} instead exploits an
environment-independent speed-distribution profile with cross-modal visual supervision,
reporting 96.8\% over a year-long deployment. Our augmentations are in the same spirit as
ARC-Fi's, and we are explicit that this is inherited rather than novel.

\subsection{Few-shot and prototype-based matching}
FewSense \cite{Yin2022FewSense} pretrains a feature extractor and fine-tunes it on a few
target samples, reaching 93.9/96.5/82.7\% five-shot accuracy on SignFi/Widar/Wiar with
collaborative multi-receiver fusion adding 29.2\%. CrossFi \cite{Zhao2024CrossFi} replaces
cosine similarity with an attention-based similarity network and a class-template
generator, reporting 91.72\% one-shot and 64.81\% zero-shot cross-domain. FewCS
\cite{Yang2023FewCS} clusters intra-domain prototypes and extracts cross-domain prototypes
in a shared space for edge--cloud online learning. WiGNN \cite{Chen2024WiGNN} models
multi-receiver CSI as a graph and averages 94.0\% cross-domain on Widar3.0. These systems
are the closest in spirit to retrieval --- matching a query against stored exemplars ---
but their memory holds \emph{class prototypes}, not \emph{spatial descriptions}, so they
cannot exploit the fact that the room is known.

\subsection{Self-supervised representation learning and foundation models}
AutoFi \cite{Yang2022AutoFi} learns from randomly captured low-quality unlabelled CSI with
a geometric self-supervised objective and transfers across tasks. WiFiAct
\cite{Lyons2024WiFiAct} pairs an environment-robust preprocessing pipeline with a Bayesian
CNN and contrastive augmentation. Context-aware predictive coding \cite{Barahimi2024CPC},
temporal--frequency masked autoencoding for localisation \cite{Liu2026SSLMAE},
cross-modal-guided MAE \cite{Liu2025CIGMAE}, CSI-MAE \cite{Jiang2026CSIMAE}, a
CIR--CSI-consistency MIMO channel foundation model \cite{Jiang2025MIMO}, JEPA-style
objectives \cite{Kim2026WiFiJEPA,Luo2026CSIJEPA}, parameter-efficient adapters
\cite{Custance2026Adapter}, source-free multi-user adaptation \cite{Radwan2026MUSHOTFi},
and a Siamese CNN--Transformer channel-charting model robust to transmitter relocation
\cite{Lee2026PCC} all build encoders without labels. The evaluation of Xu et al.\
\cite{Xu2025EvalSSL} benchmarks these objectives for CSI activity recognition, while
Jiang et al.\ \cite{Jiang2025Foundation} pretrain masked autoencoders on 1.3M samples from
14 datasets and find log-linear gains with pretraining data --- concluding that data, not
capacity, is the bottleneck --- and Gong et al.\ \cite{Gong2026WiSwiss} pre-train on
multi-source datasets. \method{} inherits this encoder-scaling result: our contribution is
what happens \emph{after} the encoder, and is complementary to pretraining scale.

\subsection{Retrieval augmentation and explicit external memory}
RAG conditions generation on retrieved non-parametric context \cite{lewis2020rag}, and
spatial variants constrain retrieval with geometric relations \cite{spatial_rag}. In
wireless, ENWAR \cite{Nazar2024Enwar} builds a RAG-empowered multimodal LLM framework for
wireless environment perception, extended in ENWAR~2.0 \cite{Nazar2026Enwar2} with adaptive
knowledge-base formation, chain-of-thought reasoning and beam tracking; Wi-Chat
\cite{wi_chat} shows an LLM can interpret CSI for activity recognition when prompted with
sensing principles. These systems put retrieval in service of language-level perception.
Wang et al.\ \cite{Wang2025GeneralizabilitySurvey} raise the open question of RAG as
on-demand Wi-Fi knowledge. To our knowledge, no prior work uses a spatially-indexed
retrieval memory as the cross-domain adaptation mechanism for a CSI sensing head, which is
the gap we target. \TODO{re-run a targeted search for "retrieval-augmented CSI sensing" and
"environment-conditioned Wi-Fi sensing" immediately before submission; this is the novelty
claim most likely to be scooped.}

% =====================================================================
\section{Method}
\label{sec:method}

\subsection{Problem formulation}
Let $x\in\mathbb{R}^{N_s\times N_t}$ denote a CSI window with $N_s$ subcarriers and $N_t$
time samples, and $y\in\{1,\dots,C\}$ its activity label. We associate each measurement
with a domain descriptor $d$ that is available \emph{at inference} without labels: a coarse
geometric signature of the space (transmitter--receiver geometry, coarse distance bins
derived from the signal statistics) optionally plus installation metadata
\DECISION{final descriptor content}. Standard cross-domain sensing trains
$f_\theta(x)\!\to\!y$ on source domains $\{d_s\}$ and is evaluated on an unseen $d_t$,
treating $d$ as nuisance. We instead model
\begin{equation}
p(y\mid x, \mathcal{M}) \;=\; \sum_{m \in \mathrm{TopK}(x,\mathcal{M})} w_m \, p_\phi\!\left(y \mid \zc(x),\, m\right),
\label{eq:main}
\end{equation}
where $\zc(x)$ is a content code, $\mathcal{M}=\{(\rho_m, e_m)\}$ is a memory of
spatial/layout entries ($\rho_m$ a geometric descriptor, $e_m$ its conditioning embedding),
and $w_m$ are retrieval weights (Eq.~\eqref{eq:main}). Adaptation to a new environment
consists of inserting an entry into $\mathcal{M}$; the parameters $\theta,\phi$ are untouched.

\subsection{Stage 1: self-supervised disentanglement}
An encoder $E_\theta$ maps $x$ to $(\zc,\zd)$ by splitting the embedding into two
sub-vectors, decoded jointly:
\begin{equation}
\hat{x} = D_\psi(\zc, \zd), \qquad
\mathcal{L}_{\mathrm{rec}} = \|x-\hat{x}\|_2^2 .
\end{equation}
Content must be predictive of the label for labelled source windows,
$\mathcal{L}_{\mathrm{cls}} = \mathrm{CE}\big(y, g_\phi(\zc)\big)$, while the domain code
must be predictable from the signal statistics but \emph{not} from the content code, and
vice versa. We enforce statistical independence with an HSIC-style penalty
\DECISION{HSIC vs adversarial classifier --- both are defensible; the adversarial variant
matches \cite{Kang2021ContextAware} and is easier to tune, HSIC is more stable}:
\begin{equation}
\mathcal{L}_{\mathrm{ind}} = \mathrm{HSIC}(\zc, \zd) \;\; \text{or} \;\;
\min_{E_\theta}\max_{h}\ \mathcal{L}_{\mathrm{adv}}(h(\zc),\zd).
\end{equation}
Finally we add physics-informed augmentations that generate views of the same event which
differ in domain but not in content --- co-located antenna responses
\cite{Xu2023ARCFi}, subcarrier-domain style randomisation \cite{Liu2024WiSR}, and
subdomain-guided perturbation \cite{Liu2024WiSGP} --- with a consistency loss
$\mathcal{L}_{\mathrm{con}}$ between the content codes of two views. The total source
objective is
\begin{equation}
\mathcal{L} = \mathcal{L}_{\mathrm{rec}} + \lambda_1 \mathcal{L}_{\mathrm{cls}}
+ \lambda_2 \mathcal{L}_{\mathrm{ind}} + \lambda_3 \mathcal{L}_{\mathrm{con}} .
\label{eq:total}
\end{equation}
\DECISION{values of $\lambda_{1..3}$, embedding split ratio, window length $N_t$,
subcarrier sampling} --- these are hyperparameters to be tuned on source domains only
(Eq.~\eqref{eq:total}), and must be reported in the final paper.

\subsection{Stage 2: spatially-indexed retrieval memory}
Each memory entry stores a geometric descriptor and the domain embedding of a reference
capture from that space,
$\rho_m = \big[\phi_{\mathrm{geom}}(x^{(m)}_{\mathrm{ref}}),\, \text{meta}_m\big]$,
$e_m = \zd(x^{(m)}_{\mathrm{ref}})$. At inference we retrieve the top-$K$ entries by
descriptor similarity, weight them by softmax over similarity, and condition the label
head on the concatenation $[\zc(x); \sum_m w_m e_m]$. Because $e_m$ is an embedding, a new
room can also be added \emph{label-free}: a short unattended capture supplies $\rho_m$ and
$e_m$. Algorithm~\ref{alg:relay} states the procedure.
\begin{equation}
w_m = \frac{\exp\!\big(s(\rho(x),\rho_m)/\tau\big)}{\sum_{k\in\mathrm{TopK}}\exp\!\big(s(\rho(x),\rho_k)/\tau\big)} .
\end{equation}

\begin{algorithm}[t]
\caption{\method{} inference (adaptation = memory insertion, no retraining)}
\label{alg:relay}
\begin{algorithmic}[1]
\Require frozen $E_\theta$, head $g_\phi$, memory $\mathcal{M}$, window $x$ from an unseen space
\State $(\zc, \zd) \gets E_\theta(x)$
\State $\rho \gets \phi_{\mathrm{geom}}(x)$
\State $S \gets \mathrm{TopK}\big(s(\rho,\rho_m)\ \forall m\big)$
\If{$\max_m s(\rho,\rho_m) < \delta$}                       \Comment{unknown space}
    \State flag \textsc{out-of-memory}; optionally insert a label-free entry from a short capture
\EndIf
\State $w_m \gets \mathrm{softmax}\big(s(\rho,\rho_m)/\tau\big)$ over $S$
\State \Return $g_\phi\big([\zc; \textstyle\sum_{m\in S} w_m e_m]\big)$
\end{algorithmic}
\end{algorithm}

\subsection{Why retrieval instead of pure invariance}
Under cross-entropy with a domain-adversarial penalty the encoder is pushed to remove
domain information. If two domains require different decision boundaries --- which is the
empirical finding behind the multi-subdomain results of \cite{Liu2024WiSGP} and the
domain-redefinition argument of \cite{Zhang2026BeyondLabels} --- invariance removes the
information needed to select the right boundary. Keeping $\zd$ and \emph{retrieving} the
matching context makes the boundary a function of retrieved state, which is exactly what
the deployment regime offers for free: an installer who knows the room. This also gives the
system an honest failure mode via the out-of-memory branch in
Algorithm~\ref{alg:relay}, in the spirit of the conformal reliability of
\cite{Li2026Conformal}.

% =====================================================================
\section{Experimental Protocol}
\label{sec:protocol}

\subsection{Claim-to-experiment map}
Table~\ref{tab:claims} fixes the mapping required of this paper: every experiment exists to
test one claim, and each row states what would falsify it. No row has been run yet; the
placeholders in Section~\ref{sec:results} fix where each outcome will be reported.

\begin{table*}[t]
\centering
\caption{Claim--experiment map. Falsification conditions are stated before the experiments
are run, to prevent post-hoc storytelling.}
\label{tab:claims}
\small
\begin{tabular}{p{0.05\textwidth}p{0.24\textwidth}p{0.29\textwidth}p{0.32\textwidth}}
\toprule
\# & Claim & Experiment / evidence & What would falsify it \\
\midrule
C1 & Disentanglement separates content from domain: $\zc$ transfers, $\zd$ does not. & Linear probe of $y$ from $\zc$ and of environment identity from $\zd$, on held-out domains; compare against a joint encoder and against an adversarially invariant encoder. & High label accuracy from $\zd$, or $\zd$ predicting environment at chance for the invariant baseline --- either one means the split is not doing what we claim. \\
C2 & Retrieval conditioning improves cross-domain accuracy over the same encoder without memory. & Memory-free ablation: identical weights, memory disabled. & No difference beyond noise across seeds. \\
C3 & Adaptation by memory insertion is competitive with adaptation that trains on target data. & Compare against DA baselines (Fidora \cite{Chen2022Fidora}, DANN-R \cite{Prasad2024DANN}, Wi-SFDAGR \cite{Yan2025WiSFDAGR}) under a \DECISION{1-shot / 5-shot} target budget. & Baselines win decisively while requiring labelled target data --- then the retrieval path has no advantage to report. \\
C4 & The out-of-memory branch detects spaces the memory does not cover. & Accuracy/calibration when querying with an environment deliberately withheld from $\mathcal{M}$. & Flag never fires, or fires on in-memory queries at a high rate. \\
C5 & Gains are not explained by encoder pretraining alone. & Match a self-supervised baseline scaled to equal pretraining data. & Baseline equals \method{} without retrieval --- then retrieval contributes nothing. \\
\bottomrule
\end{tabular}
\end{table*}

\subsection{Datasets and splits}
\DECISION{dataset selection --- pick from, and justify: Widar3.0 \cite{Zheng2019Widar3}
(gesture, the standard cross-domain benchmark: cross-location, cross-orientation,
cross-environment), SignFi / Wiar \cite{Yin2022FewSense}, CSIDA \cite{Xu2023ARCFi},
XRF55 \cite{Yan2025WiSFDAGR}. If a layout-change protocol is needed, note that Widar3.0's
domain factors already separate environment from location and orientation, which suits C1.}
Data collection for a new layout, if performed, must be reported with device counts,
sampling rate, number of subcarriers, distances, and the exact layout changes made.

\subsection{Baselines}
Include at least: (i) a supervised CNN trained on source domains only --- the naive
baseline; (ii) domain-invariant feature extraction \cite{Zheng2019Widar3}; (iii) style
randomisation / perturbation methods \cite{Liu2024WiSR,Liu2024WiSGP}; (iv) adaptation
methods \cite{Chen2022Fidora,Prasad2024DANN,Yan2025WiSFDAGR,Zheng2025AdaWiFi};
(v) few-shot prototype matching \cite{Yin2022FewSense,Zhao2024CrossFi,Yang2023FewCS};
(vi) a self-supervised encoder without retrieval \cite{Yang2022AutoFi,Barahimi2024CPC};
(vii) a reliability-oriented baseline \cite{Li2026Conformal}. Report each baseline's
compute and tag requirements in the table --- a win against a method that needed 10$\times$
the labels is not the same claim as a win under equal budget.

\subsection{Metrics and statistics}
Primary: activity accuracy and macro-F1 (higher better); for localisation variants, mean
absolute error in centimetres (lower better). Report mean $\pm$ standard deviation over
\DECISION{$\geq$ 5 seeds} and 95\% confidence intervals from \DECISION{paired bootstrap
(10k resamples)}; for paired comparisons of two methods on the same test items use
McNemar's exact test and report $p$ and the effect size, not $p$ alone. State the seed
control, the hardware, and total GPU hours. \DECISION{All of the above is design, not
result: no number in this manuscript has been measured.}

\subsection{Implementation status}
The runner, the baseline harness, the statistical analysis, and the figure generation are
specified in \texttt{HANDOFF.md} and are intended to live in
\texttt{experiments/}, \texttt{analysis/} and \texttt{figures/}. Until those exist and
have been executed, Sections~\ref{sec:results}--\ref{sec:discussion} remain scaffolds.

% =====================================================================
\section{Results}
\label{sec:results}

\RESULTNEEDED{Table 1 --- main cross-domain comparison: rows = methods, columns = domain
factor (environment, location, orientation, user), cells = accuracy/F1 mean\,$\pm$\,std
over seeds. No cell may be filled except from an executed run.}

\RESULTNEEDED{Figure 2 --- per-environment breakdown plus the memory-free ablation (C2),
and the 1-/5-shot comparison with DA baselines (C3).}

\RESULTNEEDED{Table 2 --- disentanglement diagnostics for C1: probing accuracy of the
label from $\zc$, of environment identity from $\zd$, and the independence estimate.}

\RESULTNEEDED{Table 3 --- out-of-memory detection for C4: flag rate and post-flag
accuracy on withheld environments.}

\begin{figure}[t]
\centering
\fbox{\parbox{0.92\columnwidth}{\centering\vspace{1.4em}
\textcolor{red}{\textbf{[FIGURE PLACEHOLDER]}}\\[0.4em]
\small Figure~1: system overview --- CSI window $\to$ shared encoder $\to$ ($\zc$, $\zd$)
split $\to$ top-$K$ retrieval over the geometry-tagged memory $\to$ conditioned classifier.
Draw this once the pipeline is frozen; do not reuse a schematic that depicts components the
code does not implement.\vspace{1.4em}}}
\caption{System overview. \RESULTNEEDED{replace placeholder with vector PDF
(\texttt{figures/fig1\_overview.pdf}).}}
\label{fig:overview}
\end{figure}

% =====================================================================
\section{Discussion, Limitations, and Threats to Validity}
\label{sec:discussion}

\textbf{No empirical evidence yet.} This draft reports a method and a protocol. It
contains no measured results, and no accuracy, F1, or error figure anywhere in the text
should be read as one. The only numbers that appear are those quoted from prior work and
attributed to it.

\textbf{Novelty risk.} The retrieval-for-adaptation idea sits close to two literatures:
prototype/few-shot matching \cite{Yin2022FewSense,Zhao2024CrossFi,Yang2023FewCS} and
wireless RAG \cite{Nazar2024Enwar,Nazar2026Enwar2}. The distinguishing claim --- spatially
indexed memory conditioning a sensing head, updated label-free --- must survive a fresh
search and a carefully worded related-work section. If a concurrent paper makes the same
move, the fallback contribution is the disentanglement-plus-retrieval interaction (C1, C2,
C5), which is testable independently.

\textbf{Disentanglement assumptions.} Perfect separation of content and domain is not
achievable in general; residual leakage is expected, and C1's probes are designed to
measure it rather than assert it. If $\zd$ leaks gesture information, the retrieval
conditioning becomes a soft shortcut that could inflate in-domain scores without
transferring.

\textbf{Retrieval realism.} The memory must be populated from captures that are cheap to
obtain, otherwise the approach relocates the data-collection burden rather than removing
it. The protocol must therefore report exactly what a new entry costs (duration of
unattended capture, number of devices, whether the space must be empty), and compare that
cost against the labelled data the baselines need.

\textbf{Privacy.} CSI is frequently preferred over cameras because it does not image
people \cite{Radwan2025Survey}. The geometric descriptors we store are a property of the
space, not of a person, and the memory must be specified so that it cannot be inverted into
person-identifying information. This is a design constraint, not an afterthought.

% =====================================================================
\section{Conclusion}
\label{sec:conclusion}

We propose \method{}, which treats the environment in cross-domain Wi-Fi CSI sensing as
retrievable context rather than nuisance: a self-supervised objective factorises a CSI
window into content and domain codes, and a spatially indexed memory conditions the
classifier on the space in which the measurement was taken, so that a new room is handled
by inserting an entry instead of retraining a model. We have stated the method, the
falsifiable claim--experiment mapping, and the baseline set that would establish or refute
it. The empirical work --- five claims, the runnable protocol in \texttt{HANDOFF.md} ---
remains to be executed, and no result in this manuscript has been measured. The immediate
next step is to run the source-domain training and the memory-free ablation (C1, C2) on
Widar3.0, because those two results decide whether the rest of the paper is worth writing.

\bibliographystyle{IEEEtran}
\bibliography{references}

% =====================================================================
\appendices
\section{Reference audit}
\label{app:refs}
Every citation in this paper was resolved programmatically against Crossref, arXiv, or
DBLP; no entry was written from memory. The audit table, including corrections made to the
supplied reference list (year and venue mismatches, one abstract attributed to the wrong
paper, and two entries that could not be verified at all) is in
\texttt{references\_report.md}. \TODO{one supplied entry --- Puspitasari (2014), \emph{Analisis RSSI terhadap
ketinggian perangkat Wi-Fi di lingkungan indoor} --- has no DOI and could not be
verified through Crossref, arXiv, or DBLP. It is therefore \emph{not} in the
bibliography and is marked \textbf{[CITATION NEEDED]} wherever it is used: locate its
publisher record or drop the citation.}

\end{document}
